\documentclass[letterpaper,9pt]{extarticle}
\usepackage[margin=1.15cm]{geometry}
\usepackage[spanish]{babel}
\usepackage[T1]{fontenc}
\usepackage[utf8]{inputenc}
\usepackage{lmodern,amsmath,amssymb,multicol,array}
\usepackage{enumitem}
\setlength{\parindent}{0pt}
\setlength{\parskip}{3pt}
\setlength{\columnsep}{0.65cm}
\allowdisplaybreaks
\begin{document}
\begin{center}\large\bfseries FORMULARIO PRUEBAS DE HIPÓTESIS\\PARAMÉTRICAS\end{center}
\vspace{-0.3cm}
\begin{multicols}{2}
\raggedcolumns
\textbf{$H_0:\mu=\mu_0$}
\begin{align}
Z_0&=\frac{\bar X-\mu_0}{\sigma/\sqrt n}\sim N(0,1)\\
Z_0&=\frac{\bar X-\mu_0}{s/\sqrt n}\overset{\mathrm{aprox.}}\sim N(0,1)\\
T_0&=\frac{\bar X-\mu_0}{s/\sqrt n}\sim t_{n-1}
\end{align}
\textbf{$H_0:\sigma^2=\sigma_0^2$}
\begin{equation}\chi_0^2=\frac{(n-1)S^2}{\sigma_0^2}\sim\chi^2_{n-1}\end{equation}
\textbf{$H_0:p=p_0$}
\begin{align}
Z_0&=\frac{X-np_0}{\sqrt{np_0(1-p_0)}}\sim N(0,1)\\
Z_0&=\frac{\hat p-p_0}{\sqrt{p_0(1-p_0)/n}}\sim N(0,1)\notag
\end{align}
\textbf{$H_0:\mu_1-\mu_2=\Delta_0$}\quad Grupos independientes.
\begin{align}
Z_0&=\frac{(\bar x_1-\bar x_2)-\Delta_0}{\sqrt{\sigma_1^2/n_1+\sigma_2^2/n_2}}\sim N(0,1)\\
T_0&=\frac{(\bar x_1-\bar x_2)-\Delta_0}{s_p\sqrt{1/n_1+1/n_2}}\sim t_{n_1+n_2-2}
\end{align}
Donde:
\[s_p^2=\frac{(n_1-1)s_1^2+(n_2-1)s_2^2}{n_1+n_2-2}.\]
\begin{equation}
T_0=\frac{(\bar x_1-\bar x_2)-\Delta_0}{\sqrt{s_1^2/n_1+s_2^2/n_2}}\sim t_{\nu^*}
\end{equation}
Donde:
\[\nu^*=\frac{(s_1^2/n_1+s_2^2/n_2)^2}{\dfrac{(s_1^2/n_1)^2}{n_1-1}+\dfrac{(s_2^2/n_2)^2}{n_2-1}}.\]
\textbf{$H_0:\mu_1-\mu_2=\Delta_0$}\quad Grupos pareados.
\begin{equation}
T_0=\frac{\bar d-\Delta_0}{s_d/\sqrt n}\sim t_{n-1}
\end{equation}
Donde $\bar d=\frac1n\sum_{i=1}^n d_i$ y $d_i=x_{1i}-x_{2i}$.

\columnbreak
\textbf{$H_0:p_1-p_2=\Delta_0$}
\begin{equation}
Z_0=\frac{(\hat p_1-\hat p_2)-\Delta_0}{\sqrt{\hat p\hat q(1/n_1+1/n_2)}}\sim N(0,1)
\end{equation}
Donde $\hat p=\frac{x_1+x_2}{n_1+n_2}$, $\hat q=1-\hat p$, $\hat p_1=\frac{x_1}{n_1}$ y $\hat p_2=\frac{x_2}{n_2}$.

\textbf{$H_0:\sigma_1^2=\sigma_2^2$}\quad Varianzas iguales.
\begin{equation}
F_0=\frac{s_1^2}{s_2^2}\sim F_{n_1-1,n_2-1}
\end{equation}

\textbf{Reglas de decisión}

\textbf{Regla 1.} Si el EdeP cae en la RdeR, se rechaza $H_0$. Si no cae en ella, no se rechaza $H_0$.

\textbf{Regla 2.} Si $\text{valor-}p\leq\alpha$, se rechaza $H_0$; si $\text{valor-}p>\alpha$, no se rechaza $H_0$.

\textbf{Regla 3.} El valor-$p$ es la probabilidad, bajo $H_0$, de observar un resultado al menos tan extremo como el obtenido. No es la probabilidad de que $H_0$ sea verdadera ni la probabilidad de cometer un error de tipo I.

\textbf{Códigos en R}

\scriptsize
\begin{verbatim}
# Entrada de datos
x1 <- c(7,13,6,5,5,10,8,6,7)
x2 <- c(3,7,2,3,6,2,1,0,2)

# Pruebas para una población
# z.test requiere un paquete que lo implemente
z.test(datos,mu=10,stdev=4,conf.level=0.98)
t.test(datos,mu=10,conf.level=0.98)
t.test(datos,mu=10,conf.level=0.98,
       alternative="greater")
t.test(datos,mu=10,conf.level=0.98,
       alternative="less")
prop.test(x,n,p=0.20,conf.level=0.98)

# Pruebas para dos poblaciones
t.test(x1,x2,paired=TRUE)
t.test(x1,x2,paired=FALSE,
       var.equal=TRUE,conf.level=0.98)
t.test(x1,x2,paired=FALSE,
       var.equal=FALSE,conf.level=0.98)
var.test(x1,x2)
prop.test(c(x1,x2),c(n1,n2))

paqueteDEG::test.mu(media=52,n=36,sd=8,mu0=50)
paqueteDEG::test.prop(x=45,n=100,p=0.5)                           |
paqueteDEG::test.var(sd=12,n=30,sigma20=100)`  

paqueteDEG::test.mus(mu1,n1,sd1,mu2,n2,sd2,delta0 = 0,alternative= "less",conf.level = 0.95)
paqueteDEG::test.prop(x = c(56, 42), n = c(120, 110))                                        
paqueteDEG::test.vars(sd1,n1,sd2,n2,ratio0 = 1, alternative = "greater",conf.level = 0.95) 

\end{verbatim}
\normalsize
\textbf{Tipos de pruebas}

Bilateral: $H_0:\theta=\theta_0$ frente a $H_a:\theta\ne\theta_0$ (dos colas).

Cola derecha: $H_0:\theta\leq\theta_0$ frente a $H_a:\theta>\theta_0$.

Cola izquierda: $H_0:\theta\geq\theta_0$ frente a $H_a:\theta<\theta_0$.
\end{multicols}
\end{document}
